\documentclass[11pt]{amsart}
\usepackage[T1]{fontenc}
\usepackage{lmodern}
\usepackage{amsmath,amssymb,mathtools}
\usepackage[margin=1.05in]{geometry}
\usepackage[expansion=false]{microtype}
\usepackage{tikz-cd}
\usepackage{enumitem}
\usepackage[colorlinks=true,linkcolor=blue,citecolor=blue,urlcolor=blue]{hyperref}
\newtheorem{theorem}{Theorem}[section]
\newtheorem{proposition}[theorem]{Proposition}
\newtheorem{lemma}[theorem]{Lemma}
\newtheorem{corollary}[theorem]{Corollary}
\theoremstyle{definition}
\newtheorem{definition}[theorem]{Definition}
\theoremstyle{remark}
\newtheorem{remark}[theorem]{Remark}
\DeclareMathOperator{\rank}{rank}
\DeclareMathOperator{\Lie}{Lie}
\DeclareMathOperator{\Aut}{Aut}
\DeclareMathOperator{\esssup}{ess\,sup}
\DeclareMathOperator{\Supp}{Supp}
\DeclareMathOperator{\divi}{div}
\DeclareMathOperator{\Pic}{Pic}
\newcommand{\OO}{\mathcal O}
\newcommand{\II}{\mathcal I}
\newcommand{\QQ}{\mathbb Q}
\newcommand{\CC}{\mathbb C}
\newcommand{\RR}{\mathbb R}
\newcommand{\ZZ}{\mathbb Z}
\newcommand{\PE}{\overline{\operatorname{Eff}}}
\newcommand{\ddc}{dd^c}
\title[Bounded anticanonical metrics and equivariant nonvanishing]{Bounded anticanonical metrics on klt pairs and torus quotients}
\author{OpenAI}
\date{September 26, 2026}
\hypersetup{pdftitle={Bounded anticanonical metrics on klt pairs and torus quotients},pdfauthor={OpenAI}}
\begin{document}
\begin{abstract}
Let $(W,D)$ be a projective $\QQ$-factorial klt pair with effective rational boundary. We prove that a positive Cartier multiple of $-(K_W+D)$ has a nonzero section if this divisor is nef, its pullback to a resolution admits a semipositive metric with locally bounded weights, and the resolution has nonzero structure-sheaf Euler characteristic. Consequently a smooth rationally connected projective variety with smoothly semipositive anticanonical bundle has a nonzero naturally invariant anticanonical plurisection for every algebraic torus action.
\end{abstract}
\maketitle
\tableofcontents
\section{Introduction}\label{sec:introduction}

An anticanonical nonvanishing theorem must turn positivity into a section of a specified line bundle. For a singular pair $(W,D)$, the line in question is the anti-log-canonical divisor
\[
P=-(K_W+D).
\]
Nefness says that $P$ has nonnegative degree on every curve. It does not itself construct an effective divisor linearly equivalent to a positive multiple of $P$. We prove such a conclusion when the pullback of $P$ has a semipositive metric with locally bounded weights and a resolution has nonzero structure-sheaf Euler characteristic. We then construct these hypotheses on a projective torus quotient. The quotient section returns upstairs as a naturally invariant anticanonical section.

All varieties are integral and defined over $\CC$. A pair $(W,D)$ is \emph{Kawamata log terminal}, abbreviated klt, if $W$ is normal, $K_W+D$ is $\QQ$-Cartier, and on a log resolution $\rho:\widetilde W\to W$, using compatible canonical divisors, every coefficient of $K_{\widetilde W}-\rho^*(K_W+D)$ is greater than $-1$. This includes the strict transforms of the boundary components as well as exceptional divisors. A normal variety is $\QQ$-factorial if every Weil divisor has a positive Cartier multiple.

A metric on a rational line bundle means the appropriate root of a metric on one Cartier multiple. In a local frame $e$ its weight is $\varphi=-\log |e|_h^2$; semipositivity means that $\varphi$ is plurisubharmonic, abbreviated psh. Locally bounded weights are bounded above and below on relatively compact coordinate neighborhoods. They need not be continuous.

We write $\PE(V)\subset N^1(V)_{\RR}$ for the closed cone generated by effective divisor classes in the real vector space of Cartier divisors modulo numerical equivalence. A rational Cartier divisor is \emph{pseudoeffective} if its class lies in this cone.

\begin{theorem}[Bounded-metric nonvanishing]\label{thm:klt}
Let $(W,D)$ be a projective $\QQ$-factorial klt pair with effective rational boundary, and put $P=-(K_W+D)$. Suppose that $P$ is nef. Suppose that on a smooth projective resolution $\rho:\widetilde W\to W$, the pullback $\rho^*P$ has a semipositive metric with locally bounded weights and
\[
\chi(\widetilde W,\OO_{\widetilde W})\ne0.
\]
Then $H^0(W,\OO_W(mP))\ne0$ for some integer $m>0$ such that $mP$ is Cartier.
\end{theorem}

These hypotheses persist on a common higher resolution: bounded psh weights pull back, and the structure-sheaf Euler characteristic is a birational invariant of smooth projective varieties. The effective boundary and $\QQ$-factoriality enter the cotangent-subsheaf and birational numerical arguments, respectively. The Euler hypothesis supplies the first twisted differential forms. None is suppressed when we pass to the quotient application.

If an algebraic torus $T$ acts on a smooth variety $Z$, its differential gives the \emph{natural linearization} on $-K_Z=\det T_Z$. A section is invariant when it is fixed by the resulting action on global sections.

\begin{theorem}[Invariant sections from a stable quotient]\label{thm:equivariant}
Let $Z$ be a smooth rationally connected projective variety whose anticanonical bundle has a smooth semipositive Hermitian metric. For every algebraic torus $T$ acting on $Z$, some integer $m>0$ satisfies
\[
H^0(Z,-mK_Z)^T\ne0,
\]
with respect to the natural anticanonical linearization.
\end{theorem}

The auxiliary GIT linearization used to choose a quotient and the
linearization in this conclusion have different roles. We perturb the anticanonical class by a small ample class and a small character to obtain a projective stable quotient. We then descend the \emph{original} naturally linearized anticanonical bundle. A character shift of that bundle would change the conclusion; Remark~\ref{rem:character} gives an elementary example.

Together with the finite-cover structure theorem, Theorem~\ref{thm:equivariant} also gives nonvanishing for every smooth projective variety with smoothly semipositive anticanonical bundle; see Corollary~\ref{cor:smooth}. The companion article \cite{SmoothCompanion} proves this smooth result by a finite-volume argument with prescribed poles. The present paper gives an alternative route through a bounded-metric criterion on klt pairs and a projective stable quotient. From the companion we use only its independently proved finite-cover structure and finite \'etale norm statements.

\subsection{Context and predecessors}

Two forms of nonvanishing must be distinguished. A divisor is \emph{num-effective} if it is numerically equivalent to an effective real divisor. A section of a specified positive multiple requires linear, rather than only numerical, equivalence. Lazi\'c--Matsumura--Peternell--Tsakanikas--Xie prove num-effectivity of nef anti-log-canonical divisors on projective log-canonical threefolds that are $\QQ$-factorial or have rational singularities \cite[Theorem~A]{LMPTX2023}. Their Theorem~E is particularly close to our Euler hypothesis: on a projective log-canonical pair with rational singularities, nef anti-log-canonical divisor and nonzero $\chi(\OO)$, every nef Cartier divisor of numerical dimension one is num-effective. Theorem~\ref{thm:klt} instead obtains a section of a positive multiple of the anti-log-canonical divisor, in arbitrary dimension, under its bounded-metric hypothesis.

For ordinary nonvanishing, M\"uller proves the result for projective klt pairs with nef anti-log-canonical divisor when its restriction to a general fiber of the maximal rationally connected fibration is semiample, and removes that additional premise in dimension three \cite[Theorem~A and Corollary~B]{Muller2025}. His equivariant theorem gives a naturally invariant section for projective sub-log-canonical pairs with semiample anti-log-canonical divisor and a commutative linear algebraic group \cite[Theorem~C]{Muller2025}. Thus our equivariant theorem uses a different positivity hypothesis while restricting the varieties to be smooth and rationally connected and the groups to be tori. The distinction between smooth semipositivity and semiampleness already occurs on surfaces: the classification of Chen--Filip--Sun--Tosatti--Zhang includes rational surfaces with smoothly semipositive anticanonical bundle and a smooth elliptic anticanonical curve whose normal bundle is nontorsion of degree zero \cite[Theorem~1.3]{CFSTZ2026}. Such an anticanonical bundle cannot be semiample, since its restriction to that curve has no nonzero section in any positive power.

The method of producing twisted differential forms from Euler characteristics and hard Lefschetz has an earlier antecedent in Demailly--Peternell--Schneider \cite[Theorem~2.7.3]{DPS2001}, there for pseudoeffective canonical bundles with algebraic-singularity metrics. Their Proposition~2.7.4 also passes from twisted sections to a determinant line. We use the hard Lefschetz theorem with multiplier ideals \cite[Theorem~2.1.1]{DPS2001}, together with the determinant argument developed by Lazi\'c--Peternell \cite[Lemma~4.1]{LP2018} and adapted to anticanonical geometry in \cite[Lemma~5.1]{LMPTX2023}. Our twisting exponents are already positive and unbounded, which permits the determinant step even when the anticanonical class is numerically trivial. To control that determinant, we use \cite[Theorem~4.1]{LMPTX2023}: for a projective $\QQ$-factorial log-canonical pair with effective boundary and nef anti-log-canonical divisor, the negative first Chern class of a saturated cotangent-tensor subsheaf is pseudoeffective. This is the log-pair extension of Ou's movable-slope method \cite{Ou2023}.

The next stage transfers metric positivity to the base of a fibration. Its inputs are the curvature theory of adjoint direct images of Berndtsson \cite{Berndtsson2009}, Berndtsson--P\u{a}un \cite{BerndtssonPaun2008}, and P\u{a}un--Takayama \cite{PaunTakayama2018}. We need the singular version because discrepancy rounding leaves fractional divisorial poles even when the original metric is bounded. Normalized fiber integrals then yield a bounded metric on a base line bundle. Skoda's small-exponent integrability theorem \cite{Skoda1972} and Fujino's Koll\'ar--Nadel vanishing theorem \cite[Theorem~1.3]{Fujino2018} allow one metric construction to control every nonnegative twist, including zero.

For the invariant theorem, projective geometric invariant theory \cite{MFK1994}, torus linearization \cite{Sumihiro1974,KKLV1989}, the attracting-cell geometry of Bia\l{}ynicki-Birula \cite{BB1973}, and Luna's slice theorem \cite{Luna1973} supply the quotient and its stabilizer boundary. M\"uller's use of fixed-point weights, convex geometry and ample perturbations is a close methodological predecessor \cite{Muller2025}. The additional analytic task is to bound the norms of translated lifts uniformly. Once that bound is proved, Kiselman's minimum principle \cite{Kiselman1978,DWZZ2018} gives the quotient metric, and the same bound permits extension of the resulting invariant section over the complement of the stable locus.

The smooth application belongs to the structure theory of manifolds with semipositive Ricci curvature. Yau's prescribed-Ricci theorem \cite{Yau1978} realizes a smooth semipositive anticanonical curvature form as a Ricci form. Demailly--Peternell--Schneider describe the resulting universal-cover decomposition and finite-cover Albanese fibration \cite[Structure Theorem, p.~218]{DPS1996}; Campana--Demailly--Peternell identify the remaining compact factors as rationally connected \cite[Theorem~1.4]{CDP2015}. The companion's finite-cover theorem retains the deck action in this decomposition. The need to descend through that action explains why ordinary nonvanishing on a rationally connected factor is insufficient and why the natural invariant conclusion of Theorem~\ref{thm:equivariant} is the relevant input.

\subsection{How the proofs fit together}

For Theorem~\ref{thm:klt}, Riemann--Roch and hard Lefschetz first produce twisted differentials at unbounded exponents. Their determinants give effective divisors
\[
N_i\sim_{\QQ}M+m_iP,\qquad m_i\longrightarrow\infty,
\qquad -M\text{ pseudoeffective}.
\]
We seek a rational fibration represented by maps $\pi:S\to W$ and $f:S\to Y$, with $S,Y$ smooth projective, $\pi$ birational and $f$ surjective with connected fibers, and a decomposition
\[
r\pi^*P=f^*V+H,
\]
where $r$ is a positive integer, $V$ is a line bundle on $Y$, and $H$ is effective. We require more than effectivity: above every prime divisor of $Y$, some component must be absent from $H$. That component will provide a nonvanishing local section with which to bound the metric transferred to $V$ near the boundary of the base.

To obtain this decomposition, choose the largest possible $\dim Y$ subject to $d\pi^*P-f^*A\in\PE(S)$ for some integer $d>0$ and ample line bundle $A$ on $Y$. Any further pencil controlled by $P$ and nonconstant on the general fiber would contradict this maximality. A prime divisor is \emph{horizontal} if it dominates $Y$, and \emph{vertical} otherwise. On a normal equidimensional model, the absence of such pencils forces the horizontal parts of the $N_i$ to vary affinely with $m_i$. Taking a difference and subtracting the base contribution determined by minimum coefficient-to-multiplicity ratios produces $V$ and $H$ with the required missing-component property.

The same maximality gives rank one for the adjoint direct images obtained by rounding the klt discrepancies. Their normalized integral metrics converge to a fiberwise \emph{essential supremum}, the appropriate notion for bounded measurable weights. The distinguished components of the preceding decomposition bound this metric near every base divisor, giving a bounded semipositive metric on $V$. The domination condition supplies a small singular perturbation with strict positivity in base directions and a controlled multiplier ideal. Vanishing for every nonnegative twist then gives an Euler polynomial on the base whose value at zero is positive: the rounded exceptional divisor has a section, and all higher cohomology vanishes. Thus a positive twist has a section. Multiplying by a suitable power of the section of $H$ and pushing back to $W$ gives a section of a positive multiple of $P$; the exceptional divisor disappears under pushforward. The initial Euler characteristic needed only to be nonzero.

For Theorem~\ref{thm:equivariant}, an ample perturbation gives a nonempty stable locus $U$ equal to the semistable locus. Its quotient $W=U/T$ is projective. Finite stabilizer slices give a klt boundary with coefficients $1-1/e$, where $e$ is the codimension-one ramification index, and identify $-K_Z|_U$ naturally with the pullback of $-(K_W+D)$. Convexity along algebraic arcs proves that translated lifts have no poles; compactification of their graph upgrades this to a locally uniform bound. The orbit supremum of the original norm consequently defines the bounded quotient metric. Rational connectedness supplies Euler characteristic one on a resolution. Theorem~\ref{thm:klt} produces a section. The bounded metric and compactness of $W$ bound that section globally. By the definition of the orbit-supremum norm, its pullback is bounded throughout $U$ and therefore extends across every omitted divisor in $Z$. No orbit estimate at an unstable point is required.

Section~\ref{sec:criterion} carries out the singular nonvanishing argument, and Section~\ref{sec:quotient} constructs the quotient and proves the invariant theorem. Section~\ref{sec:descent} gives the short smooth projective application using the finite-cover structure and norm results of the companion article \cite{SmoothCompanion}.

\section{A bounded-metric criterion for klt pairs}\label{sec:criterion}

Throughout this section $(W,D)$ satisfies the hypotheses of Theorem~\ref{thm:klt}, and $P=-(K_W+D)$. Divisor equalities involving line bundles mean linear equivalence; the corresponding equalities for rational divisors mean $\QQ$-linear equivalence. An assertion about an effective divisor, its support, or its coefficient is an assertion about the actual divisor. We will explicitly distinguish these uses.

\subsection{Analytic inputs and numerical comparisons}\label{subsec:inputs}

We record the analytic results in the forms used below. If $B$ is a line bundle on a compact K\"ahler $n$-fold $S$ with a semipositive singular metric $h$, the hard Lefschetz theorem \cite[Theorem~2.1.1]{DPS2001} gives a surjection
\begin{equation}\label{eq:HL}
H^0(S,\Omega_S^{n-q}\otimes B\otimes\II(h))
\longrightarrow H^q(S,K_S\otimes B\otimes\II(h)).
\end{equation}
Here $\II(h)$ tests local integrability of $|g|^2e^{-\varphi}$ in a frame of weight $\varphi$. In particular all positive powers of a bounded metric have trivial multiplier ideal.

For a projective surjective morphism $f:S\to Y$ of smooth varieties with connected fibers, the singular adjoint direct-image theorem \cite[\S3.2.1 and Theorems~3.3.4--3.3.5]{PaunTakayama2018} applies to a semipositive singular metric on $B$. When its multiplier ideal is trivial on the total space and on general fibers, its fiberwise $L^2$ metric on $f_*(K_{S/Y}+B)$ has semipositive curvature in the singular sense. On a base-change open where the direct image has rank one and its generator has finite integral, the negative logarithm of that integral is a psh local weight. The integral formula and the metric can be compared almost everywhere; we do not need continuity or base change at every point. We apply this statement separately for every integer twist, so the base-change open may depend on that integer.

Finally, the theorem of Fujino \cite[Theorem~1.3]{Fujino2018} states that if $f:S\to Y$ is a surjective morphism from a compact K\"ahler manifold to a projective variety, and a singular metric on a line bundle $B$ has curvature at least $\epsilon f^*\omega_A$ for $\epsilon>0$ and a positive curvature form of an ample line bundle $A$ on $Y$, then
\begin{equation}\label{eq:KN}
H^i\bigl(Y,R^j f_*(K_S\otimes B\otimes\II(h))\bigr)=0
\quad (i>0,\ j\ge0).
\end{equation}
Only $j=0$ will be needed.

We also need a numerical fact that explains the role of $\QQ$-factoriality. For a birational morphism $\pi:S\to W$ with $S$ smooth and $W$ normal and $\QQ$-factorial, pushforward of divisors induces a linear map on numerical divisor classes and preserves pseudoeffectivity; compare \cite[Lemma~3.4]{LMPTX2023}. Indeed, if a divisor $F$ on $S$ is numerically trivial, then $\pi^*\pi_*F-F$ is exceptional and is numerically trivial over $W$. Applying the negativity lemma to this divisor and its negative makes it zero. Intersecting $\pi^*\pi_*F=F$ with curves dominating curves of $W$ proves that $\pi_*F$ is numerically trivial. Pushforward is consequently a well-defined continuous linear map on the finite-dimensional spaces of numerical divisor classes. It sends effective divisors to effective divisors, hence preserves their closed cone. Pullback of a pseudoeffective $\QQ$-Cartier class by the birational morphisms used here also preserves pseudoeffectivity.

\subsection{Twisted differential forms and determinants}\label{subsec:determinant}

The first task is to produce an unbounded sequence of effective divisors differing by multiples of $P$ from one fixed class. We adapt the determinant argument of Lazi\'c--Peternell \cite[Lemma~4.1]{LP2018}, as used in \cite[Lemma~5.1]{LMPTX2023}. Here the original twisting exponents are positive and unbounded, so a basis containing a section of high exponent has an unbounded sum of exponents. This direct observation is sufficient in our setting; we do not import the numerical-nontriviality hypothesis of the cited lemma.

\begin{lemma}\label{lem:determinant}
Either $H^0(W,mP)\ne0$ for some positive Cartier multiple, or there are a Weil divisor $M$, strictly increasing positive integers $m_i$, and effective Weil divisors $N_i$ such that
\begin{equation}\label{eq:Ni}
N_i\sim_{\QQ} M+m_iP,\qquad -M\in\PE(W).
\end{equation}
All the $m_iP$ may be required to be Cartier.
\end{lemma}

\begin{proof}
The assertion is immediate if $W$ is a point. Suppose $\dim W=n>0$ and choose a smooth projective resolution $\rho:\widetilde W\to W$ on which the metric is bounded. Let $r_0>0$ make $r_0P$ Cartier. Riemann--Roch makes
\[
t\longmapsto\chi(\widetilde W,K_{\widetilde W}+t\rho^*(r_0P))
\]
a polynomial. Its value at $t=0$ is $(-1)^n\chi(\widetilde W,\OO_{\widetilde W})$, which is nonzero by hypothesis. It is therefore nonzero for all but finitely many positive integers. For each such integer at least one cohomology group is nonzero. Passing to an infinite subsequence fixes a cohomological degree $q$. Apply \eqref{eq:HL} with $B=t\rho^*(r_0P)$. Its multiplier ideal is trivial. We obtain a fixed $p=n-q$ and nonzero sections
\[
\sigma_i\in H^0(\widetilde W,\Omega^p_{\widetilde W}\otimes\OO_{\widetilde W}(n_i\rho^*P)),
\qquad n_i\longrightarrow\infty.
\]
If $p=0$, projection to $W$ gives the first alternative because $\rho_*\OO_{\widetilde W}=\OO_W$.

Assume $p>0$. Work over the function field $K=\CC(W)$ and choose a rational trivialization for $r_0P$. Untwisting the sections defines vectors in $\bigwedge^p\Omega^1_{K/\CC}$. Let $F_K$ be their span and $b=\dim_KF_K$. Its determinant is a one-dimensional subspace of
\[
\bigwedge^b\bigl(\bigwedge^p\Omega^1_{K/\CC}\bigr)
\subset (\Omega^1_{K/\CC})^{\otimes pb}.
\]
The indicated inclusion uses alternating tensors in characteristic zero. Intersect this line with $(\Omega_W^1)^{[\otimes pb]}$, the reflexive tensor power, and take its saturated reflexive rank-one extension. This is a sheaf $\OO_W(M)$ for a Weil divisor $M$. The determinant of any basis selected from the $\sigma_i$ gives a nonzero section of $\OO_W(M+\sum n_iP)$ on the smooth locus of $W$. It is regular at every codimension-one point: the resolution is an isomorphism there, and the original sections are holomorphic. Reflexivity extends it across codimension two. Its zero divisor is an effective Weil divisor $N$.

The sums of exponents of such bases are unbounded. Given any one of the nonzero vectors at arbitrarily high exponent, extend it to a basis using vectors from the spanning family. All exponents are positive, so the sum is at least the chosen exponent. Selecting a strictly increasing subsequence of these sums gives $m_i$ and $N_i$ in \eqref{eq:Ni}.

The pair $(W,D)$ is projective and klt, the boundary is effective, and $-(K_W+D)$ is nef. Thus \cite[Theorem~4.1]{LMPTX2023} applies to the saturated rank-one subsheaf $\OO_W(M)$ of the reflexive cotangent tensor power. It asserts that its negative first Chern class is pseudoeffective. Since $W$ is $\QQ$-factorial, this is the claimed pseudoeffectivity of the $\QQ$-Cartier divisor $-M$.
\end{proof}

\subsection{A maximal dominated fibration}\label{subsec:maximal}

Suppose we are in the second case of Lemma~\ref{lem:determinant}. Our next task is to express a multiple of a birational pullback of $P$ as the pullback of a line bundle from a suitable base plus an effective divisor. We need a further property: above each prime divisor of the base, some component must be absent from that effective divisor. This property will prevent the induced base metric from acquiring a divisorial singularity. To obtain the decomposition, we choose a base of maximal dimension whose polarization is dominated by $P$ in pseudoeffective order. Maximality will exclude pencils satisfying the corresponding domination condition that are nonconstant on a general fiber.

Consider diagrams
\begin{equation}\label{eq:diagram}
\begin{tikzcd}[column sep=large]
S \arrow[r,"f"] \arrow[d,"\pi"'] & Y\\
W &
\end{tikzcd}
\end{equation}
where $S,Y$ are smooth projective varieties, $\pi$ is birational, $f$ is surjective, and
\begin{equation}\label{eq:domination}
d\pi^*P-f^*A\in\PE(S)
\end{equation}
for an integer $d>0$ and an ample line bundle $A$ on $Y$. The map to a point is included, with the polarization omitted. Choose such a diagram with $\dim Y$ maximal.

We can require connected general fibers. Passing to the relative algebraic closure of $\CC(Y)$ in $\CC(W)$ gives a generically finite replacement of the base. The pullback of $A$ there is big. On a smooth projective model it dominates a small ample class in pseudoeffective order, so after multiplying \eqref{eq:domination} its required form persists. The same reasoning shows that smooth birational modifications of the base are harmless. In particular $\CC(Y)$ is algebraically closed in $\CC(W)$ for the resulting fibration. The finite map in the Stein factorization of the proper morphism $f$ is then birational. Since $Y$ is normal, that map is an isomorphism. Thus every fiber of $f$ is connected and $f_*\OO_S=\OO_Y$, as required by the direct-image theorem.

There is also a normal equidimensional intermediate model
\begin{equation}\label{eq:intermediate}
\begin{tikzcd}[row sep=large,column sep=large]
& S \arrow[d] \arrow[ddl,bend right=25,"f"'] \arrow[ddr,bend left=25,"\pi"] &\\
& S_0 \arrow[dl,"f_0"'] \arrow[dr,"\pi_0"] &\\
Y && W .
\end{tikzcd}
\end{equation}
Here $\pi_0$ is a birational morphism and $S$ resolves $S_0$ and the pair on $W$. To construct it, start with a projective graph of the rational map, flatten its main component by a projective modification of the base \cite{RaynaudGruson1971}, replace the base by a smooth projective modification, and normalize the main flat transform. Finite normalization preserves equidimensionality over the smooth base. Finally resolve above this model, also dominating the resolution on which the given metric is bounded. The smooth space $S$ need not be equidimensional over $Y$. We use equidimensionality only on $S_0$, where it implies that every vertical prime divisor dominates a prime divisor of $Y$.

\begin{lemma}[Excluding a new pencil]\label{lem:pencil}
In a maximal diagram \eqref{eq:diagram}, let $B$ be a line bundle on a birational model over $S$. Suppose that its class is pseudoeffectively dominated by $c\pi^*P+E'$ for some $c>0$ and a $\pi$-exceptional divisor $E'$. Then no pencil in $|B|$ is nonconstant on the general fiber of $f$. The conclusion still applies after adding to $B$ the pullback of any fixed divisor on $Y$.
\end{lemma}

\begin{proof}
Resolve the pencil, remove its fixed part, and combine its morphism to $\mathbb P^1$ with $f$. If the pencil is nonconstant on the general fiber, the image of the combined map has dimension greater than $\dim Y$. A product polarization gives a basepoint-free line bundle $F$ on the resolved source, dominated by a multiple of $\pi^*P$ plus an exceptional divisor, using \eqref{eq:domination}.

Choose a general member $F_0$ of that system. It has no $\pi$-exceptional component. Push the class comparison down to $W$, using the numerical fact in Section~\ref{subsec:inputs}; the exceptional term vanishes. The result says that $c'P-\pi_*F_0$ is pseudoeffective. Pulling back gives $c'\pi^*P-\pi^*\pi_*F_0$ pseudoeffective. Since the effective divisor $\pi_*F_0$ is $\QQ$-Cartier, its pullback contains its strict transform $F_0$ and has an effective exceptional remainder. Hence $c'\pi^*P-F_0$ is pseudoeffective. This dominates a polarization for the larger projective image. On a smooth model of that image the pulled-back polarization is big, so it dominates a small ample class; scaling and pulling back gives a diagram of the allowed form. This contradicts maximality.

For the last assertion, choose an integer $b$ such that $bA$ minus the fixed base divisor is ample. The extra pullback is then dominated by $bf^*A$, hence by $bd\pi^*P$ using \eqref{eq:domination}. Increase $c$ accordingly.
\end{proof}

\subsection{Horizontal affineness and vertical minima}\label{subsec:minima}

Pull the effective $\QQ$-Cartier divisors $N_i$ to $S_0$. We next produce the decomposition that will control a metric at the boundary of the base.

\begin{lemma}\label{lem:decomposition}
For some integer $r>0$ with $rP$ Cartier, there are a line bundle $V$ on $Y$ and an effective Cartier divisor $H$ on $S$ such that
\begin{equation}\label{eq:decomposition}
r\pi^*P=f^*V+H.
\end{equation}
For each prime divisor $G$ of $Y$, at least one component of $f^*G$ dominating $G$ has coefficient zero in $H$.
\end{lemma}

\begin{proof}
For $m_i<m_j<m_k$, set
\[
\alpha=\frac{m_k-m_j}{m_k-m_i},\qquad
\beta=\frac{m_j-m_i}{m_k-m_i}.
\]
Then $\alpha,\beta>0$, $\alpha+\beta=1$, and
\[
N_j\sim_{\QQ}\alpha N_i+\beta N_k.
\]
After clearing denominators, the two effective divisors give sections of the same line bundle. That bundle is pseudoeffectively dominated by a multiple of $P$, because $-M$ is pseudoeffective in \eqref{eq:Ni}. Lemma~\ref{lem:pencil} says that their ratio is constant on the general fiber. The ratio belongs to $\CC(Y)$, using relative algebraic closure. Its principal divisor therefore has no horizontal component on $S_0$.

It follows that the horizontal parts of $\pi_0^*N_i$ are affine functions of $m_i$ as actual divisors. Fixing the first two indices expresses them as $A_0+m_iB_0$, where the supports of $A_0,B_0$ are contained in one finite union of primes. Effectivity for an unbounded sequence of $m_i$ gives $B_0\ge0$. Consequently the difference of two members, divided by the difference of their exponents, is $\QQ$-linearly equivalent to $\pi_0^*P$ and has effective horizontal part. Clearing denominators gives a rational section of a positive Cartier multiple of $\pi_0^*P$ with no horizontal poles. Write its divisor as $T_0$.

For every base prime $G$, let
\[
v_G=\min_{F\mapsto G}
\frac{\operatorname{coeff}_F(T_0)}{\operatorname{coeff}_F(f_0^*G)}.
\]
The denominator is positive, and the minimum is over the finitely many primes dominating $G$. Only finitely many $v_G$ are nonzero. Every vertical prime of $S_0$ occurs above a base prime, by equidimensionality. Thus
\[
T_0-f_0^*\Bigl(\sum_Gv_GG\Bigr)\ge0.
\]
The difference is $\QQ$-Cartier: $T_0$ is Cartier up to its chosen multiple, and $Y$ is smooth. Multiply again to make the base divisor integral and the remainder Cartier. This gives $V$ and an effective Cartier remainder on $S_0$. Pull them to $S$ to obtain \eqref{eq:decomposition}. The component attaining the minimum has zero remainder coefficient; its strict transform on $S$ still dominates $G$ and has coefficient zero in $H$.
\end{proof}

If $Y$ is a point, \eqref{eq:decomposition} immediately gives an effective divisor linearly equivalent to $r\pi^*P$, and hence a section downstairs. From now on suppose $\dim Y>0$. We will construct a semipositive metric on $V$ from fiber integrals and use the component absent from $H$ over each base prime to prove that this metric is bounded. The integrals must come from adjoint line bundles whose fiberwise section spaces have dimension one. Discrepancy rounding supplies these line bundles, and the same maximality argument supplies their rank-one property.

\subsection{Discrepancy rounding and rank one for every twist}\label{subsec:rankone}

Use compatible canonical divisors and put
\begin{equation}\label{eq:rounding}
R=K_S-\pi^*(K_W+D),\qquad E=\lceil R\rceil,
\qquad\Xi=\lceil R\rceil-R.
\end{equation}
Choose the resolution so that the support of $R$ is simple normal crossings. Since the pair is klt, all coefficients of $R$ are greater than $-1$. Along a nonexceptional prime the coefficient is minus its coefficient in $D$, which lies in $(-1,0]$; its ceiling is zero. Thus $E$ is integral, effective and $\pi$-exceptional. The divisor $\Xi$ is effective with simple normal crossings support and all coefficients in $[0,1)$. The identity we will use is
\begin{equation}\label{eq:adjointidentity}
K_S+\pi^*P+\Xi=E.
\end{equation}

\begin{lemma}\label{lem:rankone}
For every integer $a\ge0$,
\[
\rank f_*\OO_S(E+ar\pi^*P)=1.
\]
\end{lemma}
\begin{proof}
Over the generic fiber the section of $E+aH$ gives a nonzero section, by \eqref{eq:decomposition}. This includes $a=0$. If the dimension of the generic-fiber section space were at least two, twisting by the pullback of a sufficiently ample line bundle on $Y$ would extend two independent generic sections to global sections. Their ratio would be nonconstant on the general fiber. Their line bundle is dominated by a multiple of $\pi^*P$ plus the exceptional divisor $E$, including the base twist as in Lemma~\ref{lem:pencil}. This contradicts that lemma.
\end{proof}

For $a\ge0$ define the integral line bundle
\begin{equation}\label{eq:Ba}
B_a=(ar+1)\pi^*P+\Xi=E+ar\pi^*P-K_S.
\end{equation}
The second expression proves integrality even though the two summands in the first expression can be rational divisors. The role of $\Xi$ is now precise: its fractional poles make the adjoint line integral, while keeping the multiplier ideal trivial. Initially equip it with the metric whose weights are $(ar+1)\varphi_P+\varphi_\Xi$. Here $\varphi_P$ denotes the bounded weights on $\pi^*P$, and $\varphi_\Xi$ is the divisorial weight. Its curvature is semipositive. Its multiplier ideal is trivial because the weights of $P$ are bounded and the coefficients of the simple normal crossings divisor $\Xi$ are less than one. The same is true on general fibers, after shrinking to where the boundary is relatively normal crossings.

\subsection{A bounded metric from normalized fiber integrals}\label{subsec:metric}

Our next goal is to put a bounded semipositive metric on the line bundle $V$ in \eqref{eq:decomposition}. This is where the component absent from $H$ over each base prime is used.

\begin{proposition}\label{prop:boundedbase}
The line bundle $V$ admits a singular Hermitian metric whose local weights are psh and locally bounded on all of $Y$.
\end{proposition}

\begin{proof}
Fix a dense Zariski open $Y^\circ$ over which $f$ is smooth with connected fibers, $\Xi$ is relatively simple normal crossings, and no fiber is contained in $E\cup H$. In a coordinate neighborhood in $Y^\circ$, choose frames $v$ of $V$ and $\eta$ of $-K_Y$. The section
\[
s=f^*v\,s_H
\]
of $r\pi^*P$ has divisor $H$ over that neighborhood. The section $s_Es^af^*\eta$ belongs to $K_{S/Y}+B_a$. By Lemma~\ref{lem:rankone}, it generates the direct image at a general point. On its base-change open, this generator locally identifies the adjoint direct image with $aV-K_Y$. Because the direct image has rank one, its semipositive $L^2$ metric has a psh weight equal to the negative logarithm of the squared $L^2$ norm of this generator. Dividing by $a$ gives the functions below. This is a local identification where the generator is a frame, and asserts no global invertibility at omitted base points. We obtain psh functions
\begin{equation}\label{eq:ua}
u_a(y)=-\frac1a\log\int_{S_y}t(x)^a\,d\mu_y(x),
\qquad t=|s|_{h_P^r}^2,\qquad a\ge1.
\end{equation}
The measure $\mu_y$ is the squared adjoint norm density of $s_Ef^*\eta$, using the metric on $\pi^*P+\Xi$. Locally it has a bounded positive factor, the squared modulus of a nonzero holomorphic section, and the factors $|z_j|^{-2c_j}$ with $c_j<1$. It is therefore finite on a general compact fiber and has the same null sets as a smooth volume measure. Zeros of $s_E$ and the divisorial poles occupy sets of measure zero. This description uses boundedness, not continuity, of $\varphi_P$.

For each $y_0\in Y^\circ$, choose a point on $S_{y_0}$ off $E\cup H\cup\Supp\Xi$. On a small product neighborhood for the smooth map, the holomorphic factors are bounded away from zero and the metric weights are bounded. There are a neighborhood of $y_0$, constants $c,b>0$, and open pieces in its fibers on which $t\ge c$ and whose $\mu_y$-mass is at least $b$. Thus
\[
\int_{S_y}t^a\,d\mu_y\ge b c^a,
\qquad
u_a(y)\le-\log c-\frac{\log b}{a}.
\]
These estimates extend the psh representatives across each exponent-dependent base-change exceptional set and bound them locally uniformly above on $Y^\circ$.

Take the countable intersection $Y^*$ of the full-measure sets where these integral formulas agree with the psh representatives for every positive integer $a$. This is a common full-measure subset of $Y^\circ$, not a claimed common Zariski open. For $y\in Y^*$, the elementary $L^a$-norm limit on the finite measure space $(S_y,\mu_y)$ gives
\begin{equation}\label{eq:esssup}
\lim_{a\to\infty}u_a(y)
=-\log\esssup_{S_y}|s|_{h_P^r}^2=:u(y).
\end{equation}
For completeness, if $M_y=\esssup t$, the upper bound is
$(\int t^a\,d\mu_y)^{1/a}\le M_y\mu_y(S_y)^{1/a}$.
For any $c<M_y$, the set $\{t>c\}$ has positive measure, giving the reverse limiting bound $\liminf(\int t^a\,d\mu_y)^{1/a}\ge c$. Let $c\uparrow M_y$.

The psh functions $u_a$ have locally uniformly bounded upper bounds. They cannot converge to $-\infty$ locally: on each relatively compact base neighborhood, properness, holomorphicity of $s$, and boundedness of $\varphi_P$ bound $t$ above. The measures $\mu_y(S_y)$ are bounded above there: in relative simple normal crossings charts their densities are at most $C\prod_j|z_j|^{-2c_j}$ with $c_j<1$, and properness gives a finite cover above a compact base neighborhood. Hence \eqref{eq:ua} also has a uniform lower bound on compact subsets of $Y^\circ$. Local compactness for psh functions now gives subsequential $L^1_{\mathrm{loc}}$ convergence, and \eqref{eq:esssup} identifies every subsequential limit almost everywhere. Consequently $u$ has a psh representative.

Under a frame change $v'=gv$, the function $t$ is multiplied by $|g|^2$ and $u$ changes to $u-\log|g|^2$. Changing $\eta$ contributes a term divided by $a$, which disappears in the limit. Thus $u$ defines the weights of a metric on $V$ over $Y^\circ$.

The normalized integrals have now produced positivity on the smooth part of the base. It remains to prove boundedness at the boundary; positivity alone would not suffice for the later multiplier-ideal calculation. On the inverse image of a relatively compact base coordinate neighborhood, properness and bounded weights of $h_P$ give a uniform upper bound for $|s|^2$. Formula~\eqref{eq:esssup} therefore gives a lower bound for $u$ even near a missing base fiber. For an omitted prime $G\subset Y$, choose the component $F$ above $G$ on which $H$ has coefficient zero. At a general point of $F$, the section $s$ is nonzero, the map $F\to G$ is submersive, and local transverse coordinates have the form
\[
z_1\circ f=w_1^e\,a(w),\qquad a(w)\ne0.
\]
Together with coordinates along $G$, this shows that every nearby general fiber meets a neighborhood of that point in an open set. On a smaller neighborhood $|s|_{h_P^r}^2\ge c>0$ because the holomorphic section is nonzero and the metric is bounded. The neighborhoods intersecting such a general fiber have positive $\mu_y$-measure, since that measure has the same null sets as smooth volume. No uniform positive lower bound on their measures is needed for an essential supremum. Thus the component may lie in $E$ or $\Xi$ without affecting this argument: on nearby general fibers their divisorial supports have measure zero. The essential supremum is at least $c$. Thus $u\le-\log c$ near a dense open subset of $G$.

The psh removable-singularity theorem extends $u$ across these codimension-one sets, using the upper bound just proved. What remains is an analytic subset of codimension at least two in the smooth base. Psh Hartogs extension extends $u$ there as well. Its lower bound persists by the upper-semicontinuous extension, and its psh representative has a local upper bound. The local extensions are unique and respect the frame transformations, giving the asserted bounded metric on $V$ over all of $Y$. This conclusion is local boundedness. We have used no continuity assertion for the original metric or for its essential supremum.
\end{proof}

\subsection{One small mixture and Euler interpolation}\label{subsec:interpolation}

The bounded metric on $V$ lets us represent every nonnegative integral twist $ar\pi^*P=af^*V+aH$ by a semipositive metric whose only unbounded weights come from the integral divisor $aH$. These poles will give precisely the ideal $\OO_S(-aH)$, canceling $aH$ in the adjoint bundle. To apply vanishing, we still need strict positivity in base directions. We put it on the fixed summand $\pi^*P+\Xi$, so the same construction also works for the zeroth twist.

\begin{lemma}\label{lem:mixture}
There is a metric on $\pi^*P+\Xi$ with curvature at least $\epsilon f^*\omega_A$ for some $\epsilon>0$ and with trivial multiplier ideal.
\end{lemma}
\begin{proof}
By \eqref{eq:domination}, $d\pi^*P-f^*A$ has a positive-current metric. Tensoring with a positive metric on $f^*A$ and taking a $d$th root gives a metric on $\pi^*P$ with local weights $\psi$ and curvature at least $d^{-1}f^*\omega_A$. Choose $0<\tau<1$ and mix these weights with the bounded metric:
\[
\varphi_\tau=(1-\tau)\varphi_P+\tau\psi.
\]
Its curvature is at least $(\tau/d)f^*\omega_A$. We show that $\varphi_\tau+\varphi_\Xi$ has trivial multiplier ideal for one sufficiently small $\tau$.

Choose $p>1$ so close to one that $pc_j<1$ for every coefficient $c_j$ of $\Xi$, and let $q=p/(p-1)$. On a log-resolution chart,
\[
e^{-p\varphi_\Xi}\asymp\prod_j|z_j|^{-2pc_j}
\]
is integrable. Skoda's local small-exponent integrability theorem \cite{Skoda1972} implies that $e^{-q\tau\psi}$ is integrable for sufficiently small $\tau$. A finite cover of the compact $S$ permits one choice of $\tau$ for all charts. H\"older's inequality makes $e^{-(\varphi_\Xi+\tau\psi)}$ integrable; the remaining bounded factor $e^{-(1-\tau)\varphi_P}$ does not affect integrability. This is exactly the triviality of the multiplier ideal. No statement about analytic singularities of $\psi$ is needed.
\end{proof}

\begin{proof}[Proof of Theorem~\ref{thm:klt}]
The earlier arguments already dispose of $\dim W=0$, the $p=0$ alternative of Lemma~\ref{lem:determinant}, and a zero-dimensional maximal base. Otherwise take $V,H,E,\Xi$ as constructed above. Let $h_V$ be the bounded semipositive metric on $V$ from Proposition~\ref{prop:boundedbase}. For every integer $a\ge0$, equip $B_a$ in \eqref{eq:Ba} with the tensor product of the metric from Lemma~\ref{lem:mixture}, the pullback of $h_V^{\otimes a}$, and the divisor metric of $aH$. More explicitly, if $g_H$ is a local equation of $H$, this metric has local weight
\[
\varphi_{\mathrm{mix}}+a f^*\varphi_V+a\log|g_H|^2,
\]
where $\varphi_{\mathrm{mix}}$ is a weight of the fixed metric from Lemma~\ref{lem:mixture} and $\varphi_V$ is a weight of $h_V$. Curvature is at least the same positive multiple of $f^*\omega_A$. Since $\varphi_V$ is bounded and $aH$ is integral, the multiplier ideal is exactly $\OO_S(-aH)$. Indeed, at the generic point of a component of $H$, the remaining psh weight is bounded above, so integrability forces divisibility by $g_H^a$. The quotient extends holomorphically across codimension two by normality. Conversely, after this division, the test is the trivial-ideal test in Lemma~\ref{lem:mixture}. Bounded weights do not change that ideal.

The adjoint identity becomes
\[
K_S+B_a-aH=E+af^*V.
\]
Apply \eqref{eq:KN} and the projection formula. Writing $\mathcal G=f_*\OO_S(E)$, we obtain
\begin{equation}\label{eq:allvanishing}
H^i(Y,\mathcal G\otimes\OO_Y(aV))=0
\quad\text{for every }i>0\text{ and every }a\ge0.
\end{equation}
The sheaf $\mathcal G$ is coherent and has rank one. Riemann--Roch on the smooth projective base makes
\[
Q(a)=\chi(Y,\mathcal G\otimes\OO_Y(aV))
\]
a polynomial in $a$. The effective divisor $E$ supplies a nonzero section of $\mathcal G$. At $a=0$, \eqref{eq:allvanishing} therefore gives $Q(0)=h^0(Y,\mathcal G)>0$. The polynomial is nonzero, so there exists a positive integer $a$ with $Q(a)\ne0$. Again by \eqref{eq:allvanishing}, $Q(a)=h^0(Y,\mathcal G\otimes\OO_Y(aV))$, so this space of sections is nonzero. Thus
\[
H^0(S,E+af^*V)\ne0.
\]
Multiplication by $s_H^a$ gives a nonzero section of $E+ar\pi^*P$. Normality and the effective exceptional nature of $E$ imply $\pi_*\OO_S(E)=\OO_W$: a rational function with possible poles only over codimension at least two is regular on $W$. The projection formula now gives a nonzero section of $arP$ on $W$. This is a positive Cartier multiple, as required.
\end{proof}

\section{Stable torus quotients and invariant sections}\label{sec:quotient}

Let $Z$ be smooth and projective, let $L=-K_Z$ have a smooth semipositive Hermitian metric, and let $T$ be an algebraic torus acting on $Z$. We construct a projective quotient with a klt stabilizer boundary and a bounded semipositive anti-log-canonical metric. Rational connectedness is needed only in the final application of Theorem~\ref{thm:klt}, where it gives Euler characteristic one on a resolution of the quotient.

Replace $T$ by its effective image in $\Aut(Z)$. The kernel acts trivially on the tangent bundle and hence on the natural linearization of $L$, so this does not change invariants. For the quotient construction assume that $t=\dim T>0$; the trivial-action case of Theorem~\ref{thm:equivariant} will be treated at the end. Average the weights of the original metric over the compact form $K\simeq(S^1)^t$ of $T$. This gives a smooth $K$-invariant semipositive metric $h$ on the same naturally linearized line bundle.

\subsection{The weight convention and an ample perturbation}\label{subsec:weights}

For a $T$-linearized line bundle $C$, a point $x\in Z$, and an integral one-parameter subgroup $\lambda:\CC^*\to T$, define
\begin{equation}\label{eq:weight}
w_C(x,\lambda)=\text{the weight on }C_{x_0},
\qquad x_0=\lim_{\tau\to0}\lambda(\tau)x.
\end{equation}
Projectivity gives the limit. Our convention is that a fiber vector of weight $w$ transforms by multiplication by $\tau^w$. For a very ample linearized bundle $C$, this is the maximum of the pairings with $\lambda$ of the weights of eigensections that do not vanish at $x$.

To verify the sign, let $s$ be an eigensection of weight $\alpha$ and $l$ a nonzero lift at $x$. Then the scalar of $s$ relative to $\lambda(\tau)l$ equals $\tau^{-\alpha}$ times its scalar at $x$. The smallest order among these scalars is $-\max\alpha$, and a basepoint-free collection has minimum order zero relative to a nonvanishing frame at the limit. Thus the transformed lift has order $\max\alpha$, agreeing with \eqref{eq:weight}. In this convention the Hilbert--Mumford condition is $w_C(x,\lambda)\ge0$ for semistability and strict positivity for stability, for every nonzero one-parameter subgroup; see \cite{MFK1994} and \cite[\S1.1.5]{DolgachevHu1998}.

Torus linearization and equivariant embeddings \cite{Sumihiro1974}, in the precise forms \cite[Proposition~2.4 and Corollary~2.6]{KKLV1989}, allow an equivariantly very ample line bundle $J$. Replace it by a sufficiently high power so that $J+L$ is also very ample, using the natural linearization on $L$. For each of the two very ample bundles, the nonvanishing eigensections at a point form a subset of a fixed finite list. Consequently their weight functions, and the difference
\[
w_L=w_{J+L}-w_J,
\]
have only finitely many patterns. Each pattern is homogeneous and piecewise linear on a finite rational polyhedral fan in $\Lie(T)_{\RR}$, interpreted as the real cocharacter space.

\begin{lemma}\label{lem:genericweight}
On a dense Zariski open subset of $Z$, the generic weight pattern of $L$ is strictly positive on every nonzero real cocharacter.
\end{lemma}
\begin{proof}
First fix a nonzero integral cocharacter $\lambda$. The Bia\l{}ynicki-Birula decomposition \cite{BB1973} has an open attracting stratum; the smooth attracting-piece statement and its tangent-weight description also follow from \cite[Theorem~1.5 and Corollary~7.3]{JS2018}. At its limiting fixed component, all tangent weights are nonnegative, and some are positive. Otherwise the attracting stratum would have no positive-dimensional directions transverse to a fixed component and openness would make $\lambda$ act trivially. This contradicts effectiveness. The weight of $L=\det T_Z$ is the sum of those tangent weights, so it is strictly positive on this open stratum.

The generic pattern for $J$ and $J+L$ occurs on one dense open subset, obtained by requiring every nonzero eigensection in chosen bases to be nonvanishing. This open meets the open attracting stratum for each fixed $\lambda$. Its weight function is therefore positive for every nonzero integral cocharacter. The finitely many rational linear pieces then make it positive on every nonzero real cocharacter: a zero or negative locus on one rational cone would contain a nonzero rational point. On a unit sphere this continuous homogeneous function has a positive minimum.
\end{proof}

\begin{proposition}\label{prop:stableopen}
There are a positive rational number $\delta$ and a rational character shift of arbitrarily small size for the ample rational linearized bundle $L+\delta J$ such that its stable locus $U$ is nonempty and equals its semistable locus. Every $x\in U$ satisfies
\begin{equation}\label{eq:weakweights}
w_L(x,\lambda)\ge0
\quad\text{for all integral one-parameter subgroups }\lambda.
\end{equation}
The shift is used only to define $U$; the linearization on $L$ remains natural.
\end{proposition}
\begin{proof}
The bundle $L$ is nef, so $L+\delta J$ is ample for every $\delta>0$. Subdivide the finitely many weight fans to one finite pointed rational fan. On each cone every weight pattern of $J$ and $L$ is linear.

For each pattern failing \eqref{eq:weakweights}, choose an integral ray on which its weight is negative. There are only finitely many choices. For sufficiently small $\delta>0$ and sufficiently small rational character shift, all these chosen weights remain negative. Such a pattern cannot occur at a semistable point for the perturbed linearization. On the other hand, Lemma~\ref{lem:genericweight} and its uniform positive lower bound on the unit sphere show that the generic pattern stays strictly positive under these perturbations. Thus the stable locus is nonempty.

After fixing such a rational $\delta$, choose the rational character off the finitely many hyperplanes on which a pattern has zero weight on one of the rays of the common fan. For any semistable point, weights on those rays are now strictly positive. Linearity on pointed cones makes them strictly positive on every nonzero cocharacter. The point is stable by Hilbert--Mumford. Finally clear denominators to obtain an ordinary ample linearized line bundle. This leaves the stated loci unchanged.
\end{proof}

By projective GIT the quotient
\begin{equation}\label{eq:quotient}
q:U\longrightarrow W=U/T
\end{equation}
is a geometric quotient and $W$ is projective. The action has finite stabilizers on $U$. Equality of stable and semistable loci is important here: the quotient is the complete projective GIT quotient, not merely an open subset of one.

\subsection{The stabilizer boundary and the anticanonical identity}\label{subsec:slices}

We describe the singularities of $W$ and identify the divisor to which Theorem~\ref{thm:klt} will apply.

\begin{proposition}\label{prop:boundary}
The variety $W$ is normal and $\QQ$-factorial. There is an effective rational divisor
\[
D=\sum_B\left(1-\frac1{e_B}\right)B
\]
on $W$, where $e_B$ is the generic ramification index of a finite stabilizer slice over the prime $B$, such that $(W,D)$ is klt. If $P=-(K_W+D)$, then
\begin{equation}\label{eq:canonicalquotient}
L|_U=q^*P
\end{equation}
as equivariant $\QQ$-line bundles, using the natural linearization on $L$ and the trivial action on the base.
\end{proposition}

\begin{proof}
Apply Luna's slice theorem \cite{Luna1973} in an invariant affine neighborhood coming from the ample GIT linearization. Near an orbit with finite stabilizer $G$, there is a smooth slice $S_{\mathrm{sl}}$ and, after an \emph{\'etale} base localization, the local model
\[
T\times^G S_{\mathrm{sl}}\longrightarrow S_{\mathrm{sl}}/G.
\]
The slice is smooth because $U$ is smooth and the orbit has constant dimension. Normality follows from invariants of a normal ring. Finite quotients of smooth varieties are locally $\QQ$-factorial: a Weil divisor pulls back to a Cartier divisor locally on the smooth finite cover, and the product of its translates gives a local equation for a positive multiple downstairs. The same assertion descends through the \emph{\'etale} localizations. Thus $W$ is $\QQ$-factorial.

Discard an ineffective kernel of the finite slice action when computing ramification. At a general point over a prime $B$, the finite map has ramification order $e_B$, and the ramification divisor has coefficient $e_B-1$. Thus its canonical divisor formula is
\[
K_{S_{\mathrm{sl}}}=p^*(K_W+D)
\]
in codimension one. The numbers $e_B$ are intrinsic: they are the multiplicities along components of the quotient pullback of $B$. The formulas therefore agree on overlapping charts and define one global effective rational divisor $D$. Only finitely many primes are ramified.

The usual finite-morphism discrepancy formula shows that this pair is klt. More explicitly, extend any divisorial valuation over $W$ to the function field of the slice. If the valuation has ramification index $e$, its log discrepancy upstairs is $e$ times its log discrepancy for $(W,D)$ downstairs. The upstairs variety is smooth with zero boundary, hence every log discrepancy is positive. The downstairs log discrepancies are therefore positive as well. This is the finite-morphism discrepancy criterion of \cite[Proposition~5.20]{KM1998}; the exact finite-group calculation with its ramification boundary is also given in \cite[Proposition~3.5 and equation~(3.4)]{BGLM2024}.

It remains to keep track of the linearization in \eqref{eq:canonicalquotient}. On $U$ the infinitesimal action gives an injective subbundle
\[
\Lie(T)\otimes\OO_U\ \longrightarrow\ T_U.
\]
Its rank is $\dim T$ because all stabilizers are finite. Choose a dense smooth open $W^\circ\subset W$ away from the branch boundary over which $q$ is a submersion, and put $U^\circ=q^{-1}(W^\circ)$. There the tangent sequence is
\[
0\longrightarrow\Lie(T)\otimes\OO_{U^\circ}
\longrightarrow T_{U^\circ}
\longrightarrow q^*T_{W^\circ}\longrightarrow0.
\]
The adjoint representation of a torus is trivial, so a nonzero element of $\det\Lie(T)$ gives an equivariant trivialization of the orbit determinant. Taking determinants of this sequence therefore identifies $L|_{U^\circ}$ with $q^*(-K_{W^\circ})$ equivariantly. On each slice the orders of this rational identification are exactly those in the ramification calculation above, which introduces the boundary $D$. No character twist occurs.

The quotient has constant fiber dimension $\dim T$. A prime divisor in $U$ cannot map into codimension at least two in $W$: the dimension bound for such a preimage would be at most $\dim U-2$. Thus the preceding slice and codimension-one calculation accounts for every prime divisor upstairs. It gives the equivariant rational-line-bundle identity \eqref{eq:canonicalquotient}. Equivalently, choose finitely many slice charts covering the projective quotient. Their finite stabilizer groups have bounded exponent, and a common multiple also clears the Cartier indices. At that multiple we obtain an actual isomorphism of line bundles whose action on the pullback from $W$ is trivial along each fiber. This is an equivariant rational-line-bundle identity, not only a numerical equality.
\end{proof}

\subsection{A uniform bound on translated lifts}\label{subsec:orbitbound}

The quotient is now algebraically suitable. To apply Theorem~\ref{thm:klt} and later extend its section, we need a locally bounded quotient metric. We prove the required analytic estimate instead of inferring it from the existence of the GIT quotient.

\begin{lemma}[Orbit bound]\label{lem:orbitbound}
Let $x_0\in U$ and let $l(x)$ be a nonvanishing algebraic local frame of $L$ near $x_0$. Then
\begin{equation}\label{eq:orbitbound}
\sup_{t\in T}|t\cdot l(x)|_h
\end{equation}
is bounded above on a neighborhood of $x_0$. The same holds for holomorphic local frames and positive tensor powers of $L$.
\end{lemma}

\begin{proof}
We first prove that translated lifts have no poles along algebraic arcs starting in $U$. A compactification of the action graph will then turn this statement about orders into a locally uniform bound. Let $s\mapsto x(s)$ tend to $x_0\in U$, let $l(s)$ be a nonzero holomorphic lift with nonzero value at $s=0$, and let $t(s)\in T$ be meromorphic at zero. Since a torus is a product of multiplicative groups, write
\[
t(s)=\lambda(s)u(s)
\]
with an integral cocharacter $\lambda$ and a holomorphic map $u$ into $T$ whose value at zero is invertible. Replace $x(s),l(s)$ by $u(s)x(s),u(s)l(s)$. Rename their new initial point $u(0)x_0$ as $x_0$. It is still in $U$, which is invariant, so the weak-weight inequality remains available. It suffices to consider $t(s)=\lambda(s)$.

To compare the original lift with $\lambda(s)l(s)$, vary the rate at which its cocharacter parameter tends to zero. For rational $r\in[0,1]$, let $b(r)$ be the order of $\lambda(s^r)l(s)$ relative to a nonvanishing local frame at the limiting base point. Positive order means a zero of the lift and negative order means a pole. Take a common ramified parameter to define this order and divide by its degree. Projectivity gives the limiting base point, and algebraicity gives a finite integral order before normalization. Smoothness and positivity of the metric imply
\begin{equation}\label{eq:arcorder}
b(r)=\lim_{s\downarrow0}
\frac{-\log|\lambda(s^r)l(s)|_h^2}{-\log|s|^2}.
\end{equation}
For each fixed real $s>0$, the numerator is convex as a function of $r$. Indeed the pullback metric along the holomorphic orbit map from $\CC^*$ has psh weight, and invariance under the compact circle makes that weight convex in the real logarithmic coordinate. Composition with $r\mapsto r\log s$ preserves convexity. The positive denominator in \eqref{eq:arcorder} preserves it too. Passing to limits shows that $b$ is convex on rational points, and $b(0)=0$.

We claim that for sufficiently small positive rational $r$,
\begin{equation}\label{eq:initialslope}
b(r)=r w_L(x_0,\lambda).
\end{equation}
To prove the claim without assuming $L$ ample, first use the very ample bundles $J$ and $J+L$ from Section~\ref{subsec:weights}. Choose a nonvanishing local lift to $J$ along the arc and tensor it with $l(s)$ to obtain the lift to $J+L$. For either very ample bundle $C$, choose an eigenbasis of sections, with weights $\alpha_j$. Let $\nu_j\ge0$ be their orders relative to the initial lift; set $\nu_j=+\infty$ for an identically vanishing section. Relative to the transformed lift their orders are
\[
\nu_j-r\alpha_j.
\]
Because the sections generate $C$, the order of the transformed lift is
\[
-\min_j(\nu_j-r\alpha_j).
\]
For small enough $r>0$, every minimizer has $\nu_j=0$. The order is therefore $r\max_{\nu_j=0}\alpha_j=r w_C(x_0,\lambda)$. Subtract the formula for $J$ from that for $J+L$ to obtain \eqref{eq:initialslope}.

For such an $r$, convexity and $b(0)=0$ give
\[
b(r)\le r b(1),\qquad
b(1)\ge\frac{b(r)}r=w_L(x_0,\lambda)\ge0,
\]
where the last inequality is \eqref{eq:weakweights}. Thus a transformed lift has no pole along any such algebraic arc.

We now prove uniformity in $x$. Suppose \eqref{eq:orbitbound} were unbounded as $x\to x_0$. Compactify $T$ projectively, and compactify the total space of $L$ as $\mathbb P(\OO_Z\oplus L)$ with its infinity divisor. Consider the graph of the algebraic map $(t,x)\mapsto t\cdot l(x)$, retaining $x$ as a separate parameter, and take its closure in the resulting proper space over the parameter neighborhood. Since $Z$ is projective and $h$ is smooth, a sequence of unbounded norms gives a point of this closure at fiber infinity above $x_0$. Algebraic curve selection through that point, with generic point in the graph, gives after normalization an arc $x(s)\to x_0$ and a meromorphic torus arc $t(s)$ for which $t(s)l(x(s))$ has a pole. This contradicts the preceding paragraph. Hence the bound is locally uniform.

A holomorphic frame differs from an algebraic frame by a nowhere vanishing holomorphic function, bounded above and below on a smaller neighborhood. Positive powers preserve the estimate. This proves the remaining assertions.
\end{proof}

\subsection{The orbit-supremum metric}\label{subsec:minimum}

Fix a divisible positive integer $N$ for which \eqref{eq:canonicalquotient} is an isomorphism of line bundles. For $y\in W$ and a vector $v\in (NP)_y$, pull $v$ back to the orbit $q^{-1}(y)$ and define
\begin{equation}\label{eq:supnorm}
|v|_{h_{NP}}=\sup_{x\in q^{-1}(y)}|q^*v(x)|_{h^N}.
\end{equation}
The definition is independent of a representative of the orbit, since the isomorphism is equivariant. It scales by $|c|$ on replacing $v$ by $cv$. Finite stabilizer characters have been killed by the choice of $N$.

\begin{proposition}\label{prop:quotientmetric}
The norm \eqref{eq:supnorm} defines a metric whose $N$th root on $P$ has locally bounded psh weights on $W$, and whose pullback to every smooth resolution also has locally bounded psh weights. In particular $P$ is nef.
\end{proposition}

\begin{proof}
Consider a finite slice chart as in Proposition~\ref{prop:boundary}. Nearby quotient points have representatives on the slice near a fixed representative. Lemma~\ref{lem:orbitbound} gives a local upper bound for the norm of the pullback of a frame of $NP$. The supremum also includes the norm at the representative itself. The latter is bounded away from zero in a small chart, since the original metric is smooth and the frame is nonvanishing. Thus the supremum is finite and bounded above and below in frames. Its negative logarithm is consequently locally bounded.

To prove psh positivity, work over a small smooth slice coordinate ball $B$ and the universal cover $\CC^t\to T$. Let $e$ be the pulled-back local frame of $NP$ and set
\[
\Phi(z,\zeta)=-\log\bigl|\exp(\zeta)\cdot e(z)\bigr|_{h^N}^2
\qquad(z\in B,\ \zeta\in\CC^t).
\]
The action map is holomorphic and the translated frame is holomorphic and nonvanishing, so $\Phi$ is jointly psh. The metric is invariant under the compact form of $T$, hence $\Phi$ is independent of $\operatorname{Im}\zeta$. The domain $B\times\CC^t$ is pseudoconvex and has connected tube fibers. Kiselman's minimum principle \cite{Kiselman1978}, in the form of \cite[Theorem~1.1]{DWZZ2018}, therefore makes
\[
\inf_{\zeta\in\CC^t}\Phi(z,\zeta)
\]
psh or identically $-\infty$. It is the weight of \eqref{eq:supnorm}, and the bounds already proved exclude the latter possibility. Notice that the minimum principle supplies positivity after boundedness has been established; it does not supply the orbit bound.

The resulting slice weights are invariant under the finite slice group and have the correct transformations for the pullback of $NP$. They descend to locally bounded upper-semicontinuous functions on the finite quotient. To check positivity, take any holomorphic disc in that quotient. Normalize a component of its finite pullback to the slice that dominates the disc. After a finite ramified disc cover it has a holomorphic lift to the slice. The pulled-back weight is subharmonic there, hence subharmonic downstairs off the branch points; boundedness extends it across those points. This proves the disc test for the descended weight. In particular, after composition with any holomorphic disc in a smooth resolution, the same argument gives a subharmonic function. The pulled-back weights on the resolution are therefore psh. Equivalently, extend from the regular locus using the local upper bounds and \cite[Chapter~I, Theorem~5.24]{Demailly2012}, which is applied on the smooth resolution. Bounds above and below persist. Taking the $N$th root gives the claimed metric on $P$.

Finally let $C$ be an irreducible curve on a resolution. A local bounded psh weight restricts to a psh function on the normalization of $C$, rather than being identically $-\infty$. The resulting positive curvature measure has total mass equal to the degree of the restricted line bundle. This degree is nonnegative. Thus the pullback of $P$ has nonnegative intersection with every curve, so is nef on the projective resolution. Curves downstairs are dominated by curves upstairs; the projection formula gives nefness of $P$ on $W$.
\end{proof}

\subsection{The Euler premise and extension over the complement}\label{subsec:extension}

\begin{proof}[Proof of Theorem~\ref{thm:equivariant}]
Now assume that $Z$ is rationally connected. If the effective image of $T$ is trivial, every section is naturally invariant. Theorem~\ref{thm:klt} applied to $(Z,0)$ gives such a nonzero section, since $L$ is nef and has a smooth semipositive metric and $\chi(Z,\OO_Z)=1$.

Otherwise take the projective quotient $(W,D)$ of Proposition~\ref{prop:boundary}. It is $\QQ$-factorial and klt with effective rational boundary. Proposition~\ref{prop:quotientmetric} proves the nefness and bounded-metric hypotheses for $P=-(K_W+D)$.

Let $\widetilde W\to W$ be a smooth projective resolution. The rational map $Z\dashrightarrow\widetilde W$ is dominant. Resolve its indeterminacy to a morphism $\widetilde Z\to\widetilde W$. Rational connectedness is preserved by smooth projective birational modification and by a dominant rational map to a smooth projective target; see \cite{Kollar1996}. Alternatively, pullback of differential forms under this dominant morphism is injective and $H^0(\widetilde Z,\Omega^p_{\widetilde Z})=0$ for $p>0$. Hence $H^0(\widetilde W,\Omega^p_{\widetilde W})=0$ for $p>0$, and Hodge symmetry yields
\[
\chi(\widetilde W,\OO_{\widetilde W})=1.
\]
If $W$ is a point, this identity and the same conclusion hold directly.

Theorem~\ref{thm:klt} provides a nonzero section $\sigma$ of a positive Cartier multiple $mP$. Take a positive power so that $m$ is also divisible by the integer used in \eqref{eq:canonicalquotient}. Its pullback is a nonzero invariant section $\widetilde\sigma$ of $mL$ on $U$. The bounded quotient metric makes $|\sigma|$ bounded on the compact $W$: in finitely many relatively compact charts this follows from bounded holomorphic coefficients and the bounded frame norms. Definition~\eqref{eq:supnorm} implies
\begin{equation}\label{eq:extensionbound}
|\widetilde\sigma(x)|_{h^m}
\le |\sigma(q(x))|_{h_{mP}}\le C
\quad (x\in U).
\end{equation}

Let $x$ be any point of $Z\setminus U$ and choose a holomorphic nonvanishing frame of $mL$ on a small coordinate neighborhood. Smoothness of $h$ gives upper and lower positive bounds for the norm of that frame. The coefficient of $\widetilde\sigma$ is therefore a bounded holomorphic function off the algebraic set $Z\setminus U$. The removable-singularity theorem extends it holomorphically across that set \cite[Chapter~I, Corollary~5.25]{Demailly2012}. In particular the estimate applies at the generic point of every omitted divisor; no exception for an unstable divisor is made. The same argument, or normality, extends across the remaining codimension-two locus. Unique local extensions glue to a global section of $mL$.

For any $t\in T$, the translated global section agrees with the original on the dense open set $U$, so it agrees everywhere. The section remains nonzero since it was nonzero on $U$. This proves the invariant nonvanishing assertion.
\end{proof}

\begin{remark}\label{rem:character}
Natural linearization cannot be dropped even on $\mathbb P^1$. For the action $t\cdot[z_0:z_1]=[tz_0:z_1]$, the weights on $H^0(\mathbb P^1,-mK_{\mathbb P^1})$ under the natural linearization run from $-m$ to $m$. Tensor the linearization on $-K_{\mathbb P^1}$ by the character $t\mapsto t^2$. The weights in its $m$th power then run from $m$ to $3m$, so every positive power has zero invariant subspace. This is why the auxiliary character used to choose the stable open was not transferred to $L$.
\end{remark}

\section{A smooth projective application}\label{sec:descent}

A section on a compact factor of the universal cover descends only if it respects the deck action. The invariant conclusion of Theorem~\ref{thm:equivariant} supplies precisely this compatibility with the residual torus monodromy. We obtain another proof of the smooth nonvanishing result established in \cite[Theorem ``Anticanonical nonvanishing'']{SmoothCompanion} by a finite-volume argument. Here we use that article's finite-cover structure theorem and norm lemma; the invariant section comes from the bounded-metric criterion and stable quotient constructed above.

\begin{corollary}\label{cor:smooth}
Let $X$ be a smooth connected projective complex variety with smoothly semipositive anticanonical bundle. Then $H^0(X,-mK_X)\ne0$ for some integer $m>0$.
\end{corollary}

\begin{proof}
Use the finite-cover structure theorem of \cite[Theorem ``Finite cover with compact torus monodromy'']{SmoothCompanion}, in its form retaining the deck action. There is a finite \'etale cover $X'\to X$ whose universal cover is
\[
\widetilde X=\CC^q\times B\times Z.
\]
The compact factors are projective; $K_B$ is trivial, and $Z$ is rationally connected with smoothly semipositive anticanonical bundle. Either factor may be a point. The deck group acts by translations on $\CC^q$, trivially on $B$, and on $Z$ through a compact torus contained in an algebraic torus $T$. The Euclidean and $B$ factors have invariant nowhere zero canonical frames.

By Theorem~\ref{thm:equivariant}, choose a nonzero $T$-invariant section $s_Z$ of $-mK_Z$ for some $m>0$. If $Z$ is a point, use the constant section. Tensor $s_Z$ with the $m$th powers of the dual canonical frames on the other two factors. This yields a nonzero section of $-mK_{\widetilde X}$ invariant under the deck group of $\widetilde X\to X'$. Natural linearization is crucial here: the action on the anticanonical product is the differential action used in Theorem~\ref{thm:equivariant}. The section descends to $X'$, and projectivity makes it algebraic. If $d$ is the degree of $X'\to X$, the finite \'etale norm \cite[Lemma ``Finite \'etale norm'']{SmoothCompanion} gives a nonzero section of $-mdK_X$.
\end{proof}

The geometric input in this proof rests on Yau's prescribed-Ricci theorem \cite{Yau1978} and the structure results of Demailly--Peternell--Schneider \cite{DPS1996} and Campana--Demailly--Peternell \cite[Theorem~1.4]{CDP2015}. A finite cover need not be a product: the residual torus monodromy is the reason that an invariant section, rather than an arbitrary section on $Z$, is needed.

The positive multiple in Corollary~\ref{cor:smooth} cannot in general be replaced by the first power. For an Enriques surface $S$, $K_S$ is a nontrivial line bundle of order two and $H^0(S,-K_S)=0$ \cite{Dolgachev2016}. A flat metric on $-K_S$ and the Fubini--Study metric on $\OO_{\mathbb P^1}(2)$ make the anticanonical bundle of $S\times\mathbb P^1$ smoothly semipositive. The product formula for sections shows that it has no first-power section, while its second power has sections. Its restriction to a $\mathbb P^1$ fiber has degree two, so this example is not explained by torsion of the anticanonical bundle itself.

\bibliographystyle{amsplain}
\bibliography{references}
\end{document}
